% This file is part of Oaklisp.
%
% This program is free software; you can redistribute it and/or modify
% it under the terms of the GNU General Public License as published by
% the Free Software Foundation; either version 2 of the License, or
% (at your option) any later version.
%
% This program is distributed in the hope that it will be useful,
% but WITHOUT ANY WARRANTY; without even the implied warranty of
% MERCHANTABILITY or FITNESS FOR A PARTICULAR PURPOSE.  See the
% GNU General Public License for more details.
%
% The GNU GPL is available at http://www.gnu.org/licenses/gpl.html
% or from the Free Software Foundation, 59 Temple Place - Suite 330,
% Boston, MA 02111-1307, USA


\chapter{The Compiler}

\section{File Types}

There are a number of different kinds of object files, distinguished
by extension.

\begin{center}
\begin{tabular}{@{\hspace{1.5em}}ll}\toprule
\multicolumn{1}{l}{\textbf{extension}} & \multicolumn{1}{l}{\textbf{file type}} \\\cmidrule(r){1-1}\cmidrule{2-2}
\tt .oak  & Oaklisp source file \\
\tt .omac & Macroexpanded Oaklisp source file \\
\tt .ou   & Assembly file, not peephole optimized \\
\tt .oc   & Assembly file, peephole optimized \\
\tt .oa   & Assembled object file \\\bottomrule
\end{tabular}
\end{center}

\gv{compiler-from-extension}
\doc{The extension of the input files the compiler will read.
Default \df{".oak"}.  This variable is in the compiler locale.}

\gv{compiler-to-extension}
\doc{The extension the the output files the compiler will produce.
Default \df{".oa"}.  This variable is in the compiler locale.}

\gv{compiler-noisiness}
\doc{The amount of noise the compiler should produce; zero for none, 1
for a little, and 2 for a lot.  Default value is 1, but the
\df{oakliszt} batch file compiler sets it to zero.  This variable is
in the compiler locale.}


\section{Object File Formats}

A \texttt{.oa} file is read with \df{read} and holds one form, so it
is the external representation of a list and nothing more.  A header
line may come first:

\begin{flushleft}\tt
;oaklisp-bytecode word-size=32 instructions-per-ref=2
\end{flushleft}

naming the architecture the file was assembled for.  The loader and
both cold linkers check the packing against their own; the word size
matters only because integer constants were range checked against
it.  A natively compiled file carries a second header line,
\texttt{;oaklisp-native \emph{name}\texttt{.so}}, naming the shared
object to load alongside it.

The form itself is

\begin{flushleft}\tt
(\emph{symbol-vector} (\emph{block} \ldots))
\end{flushleft}

The symbols a file mentions are collected into one vector at the
front and referred to by index, which keeps a file that names the
same symbol a hundred times from writing it a hundred times.
\df{make-oa-list} and \df{make-oaf-list} in \texttt{load-oaf.oak}
convert between that and the form the loader and the assembler work
in, where each reference names its symbol outright.

Each \emph{block} is

\begin{flushleft}\tt
(\emph{patches} \emph{opcodes})
\end{flushleft}

\emph{Opcodes} is a flat list of 16-bit instruction numbers, of even
length, which the loader packs
\df{\%\%instructions-per-ref} to a reference.  \emph{Patches} is a
flat list of triples --- \df{triplify!} cuts it into threes without
consing --- each of them a kind, an offset into the block, and a
value.  A kind of 5 or more means the value is an index into the
symbol vector and the real kind is five less.  The kinds are 0 for a
global variable, 1 for a reference to another block, and 2 for a
constant.  A constant may be a symbol, \df{()}, \df{\#t}, a number, a
character, a pair or a string.

The top-level block of a file is distinguished by the chain the cold
linker builds through it: see Chapter~\ref{sec:boot} and
\texttt{tool.oak}.

The intermediate extensions are the same data at earlier stages.
A \texttt{.omac} file is macroexpanded source, still a list of forms.
A \texttt{.ou} file is the instruction list the code generator
produced, and a \texttt{.oc} file is that list after the peephole
optimizer.  Both are written only when asked for; the ordinary path
from \texttt{.oak} to \texttt{.oa} holds them in memory.
\df{include-stage?} in \texttt{file-compiler.oak} decides which
stages a given run passes through, by numbering the extensions in
the order above.

\section{Compiler Internals}

\subsection{The Representation}

A form is read into a tree of AST nodes (\texttt{mac-compiler-nodes.oak}),
walked by a series of passes each of which fills in more of the slots,
and finally turned into a list of instructions.  Every node knows
the contour it sits in, which is how a variable reference finds out
where its variable lives.

The node types are \df{ast-constant-node}, \df{ast-variable-node},
\df{ast-if-node}, \df{ast-set-node}, \df{ast-block-node},
\df{ast-combination-node} for a call, \df{ast-make-locative-node},
\df{ast-catch-node}, and the two contours: \df{ast-method-node} and
\df{ast-labels-node}.  A contour binds names, and carries the lists
that say which of them can live on the stack and which have been
closed over and must be moved to the heap.

A variable reference is classified by \df{var-type} into
\df{stack} (a local in a frame that is still there), \df{ivar} (a
slot of the object the method is running on), \df{evar} (closed
over, and so reached through the environment) and \df{global}.

The AST walkers are \df{walk-ast} and \df{map-ast!}, with
\df{-with-arg} variants that pass a second argument along.  A pass is
written as a method per node type that does its own work and then
calls one of these, so no pass needs to know the shape of a node it
does not care about.  By convention a pass whose name ends in
\texttt{-v} may replace nodes and so its value matters; one without
the suffix only changes nodes in place.

\subsection{The Passes}

\df{compile} (\texttt{mac-compiler2.oak}) runs them in this order.
\df{expand-compile} is the same preceded by macroexpansion, and is
what the file compiler and the interpreter call.

\begin{itemize}
\item \df{code->ast} reads the form into a tree.

\item \emph{Annotation}: \df{mark-inline-lambdas} finds the
applications of a literal lambda that can be compiled as a binding
rather than a call; \df{mark-tail-nodes} finds the calls in tail
position relative to a \df{labels}; and
\df{top-level-resolve-variables} classifies every variable
reference.

\item \df{rewrite-labels-v} rewrites a \df{labels} that cannot be
compiled with jumps --- one whose labels are referred to anywhere but
as a call in tail position --- into \df{let} and \df{set\protect\bang}.
Doing so can expose more work, so annotation and this pass alternate
until the count of labels stops changing.

\item \df{fold-constants-v} replaces what can be worked out now with
the answer: a call to a foldable operation on constants, a reference
to a frozen variable.  This is also where a method's instance
variable map is guessed, since folding is what turns the type
expression into the type itself.

\item \df{insert-method-maker-v} replaces a method node with the code
that will build and install the method at run time, which is a call
to \df{\%install-method-with-env} or one of its relatives.

\item \df{heapify-variables-v} moves the variables that are closed
over off the stack: \df{determine-heap-vars} finds them by looking
for \df{\%make-locative}, \df{indirect-to-heap-v} rewrites the
references to go through the locative, \df{fold-contents-v} undoes
that again where the reference was itself a \df{make-locative} or an
assignment, and \df{insert-heap-prologue} emits the
\df{\%make-cell} calls at the head of the method.

\item \df{stash-stack-map} records in each node what the value stack
looks like where it runs, which is what a continuation needs in
order to say where the live values are.
\end{itemize}

\df{gen-code} (\texttt{mac-code.oak}) then walks the finished tree.
Each node generates its value and then its continuation, which is
the symbol \df{tail}, the symbol \df{next}, or a label;
\df{nguys-to-pop} says how much of the frame is dead,
\df{pop-args} emits the popping, and
\df{primitive-continuation} puts the two together and either falls
through, jumps, or returns.  \df{peephole-optimize}
(\texttt{peephole.oak}, from the rules compiled by
\texttt{crunch.oak}) cleans up after it, and \df{assemble}
(\texttt{assembler.oak}) resolves the labels and lays the
instructions out, padding inline arguments to a reference boundary
according to \dffl{target-architecture}.

\subsection{Native Compilation}

\df{native-compile-file} takes the same path as far as the
peephole-optimized instruction lists and then translates each block
into a C function instead of assembling it.  The \texttt{.oa} file it
writes is ordinary bytecode whose blocks are stubs, so a world
without the shared object can still load it and run the escapes.
This is described in \texttt{CompilerPlan.md} and in the comments at
the top of \texttt{src/emulator/oak-native.h}.
